Fix integers $n\ge 1$ and $m\in\mathbb Z$. To remove an endpoint ambiguity in the source text, I will work with the following convention:

- For an integer $L\ge 1$, write $(m,m+L]_{\mathbb Z}:=\{m+1,m+2,\dots,m+L\}$.

Define $f(n,m)$ to be the minimum $L\in\mathbb N$ such that there exist \emph{distinct} integers
\[
 a_1,\dots,a_n\in (m,m+L]_{\mathbb Z}
\]
with
\[
 k\mid a_k\quad\text{for every }k\in\{1,2,\dots,n\}.
\]
(If one uses the open interval $(m,m+L)\cap\mathbb Z$ instead, the value changes by at most $1$; none of the asymptotic questions in the problem are affected.)

The problem asks for proofs of the following two asymptotic statements as $n\to\infty$:

(A) \emph{Near-linear worst-case length:}
\[
\max_{m\in\mathbb Z} f(n,m) \le n^{1+o(1)}.
\]
Equivalently: for every $\varepsilon>0$ there exists $n_0(\varepsilon)$ such that for all $n\ge n_0(\varepsilon)$,
\[
\max_m f(n,m) \le n^{1+\varepsilon}.
\]

(B) \emph{Unbounded offset from the diagonal:}
\[
\max_{m\in\mathbb Z}\bigl(f(n,m)-f(n,n)\bigr)\to\infty.
\]

Stress points / edge cases: $n=1,2$ (definition sanity), dependence on endpoint convention (changes by $O(1)$), and the possibility that the maximizing $m$ is highly structured (e.g. $m$ divisible by many small integers).
